% Controlled Unicode-compatible fixture derived from the supplied HCMUT main.tex.
% The untouched supplied template remains the historical source reference.
\documentclass[a4paper]{article}
\usepackage{fontspec}
\usepackage{polyglossia}
\setdefaultlanguage{vietnamese}
\setmainfont{Times New Roman}
\usepackage{geometry,graphicx,fancyhdr,lastpage,hyperref,caption,booktabs,longtable,array,tabularx}
\usepackage{amsmath,amssymb,amsthm}
\usepackage{xcolor,listings,tikz}
\usetikzlibrary{calc}
\geometry{a4paper,total={160mm,240mm},left=30mm,top=22mm,right=20mm,bottom=20mm}
\setlength{\headheight}{40pt}
\pagestyle{fancy}
\fancyhead{}
\fancyhead[L]{
  \begin{tabular}{rl}
    \begin{picture}(25,15)(0,0)
      \put(0,-8){\includegraphics[width=10mm,height=10mm]{Images/hcmut.png}}
    \end{picture}&
    \begin{tabular}{l}
      \textbf{Trường Đại Học Bách Khoa - Đại học Quốc gia TP.HCM}\\
      \textbf{KHOA KHOA HỌC VÀ KỸ THUẬT MÁY TÍNH}
    \end{tabular}
  \end{tabular}
}
\fancyfoot{}
\fancyfoot[L]{\scriptsize Báo cáo Đồ án môn học - HK1 2026 - 2027}
\fancyfoot[R]{\scriptsize Trang {\thepage}/\pageref{LastPage}}
\renewcommand{\footruleskip}{1mm}
\renewcommand{\headrulewidth}{0.3pt}
\renewcommand{\footrulewidth}{0.3pt}
\addtocontents{toc}{\protect\thispagestyle{empty}}
\addtocontents{lof}{\protect\thispagestyle{empty}}
\addtocontents{lot}{\protect\thispagestyle{empty}}
\setcounter{secnumdepth}{4}
\setcounter{tocdepth}{3}
\makeatletter
\newcounter{subsubsubsection}[subsubsection]
\renewcommand\thesubsubsubsection{\thesubsubsection.\@alph\c@subsubsubsection}
\newcommand\subsubsubsection{\@startsection{subsubsubsection}{4}{\z@}{-3.25ex\@plus -1ex \@minus -.2ex}{1.5ex \@plus .2ex}{\normalfont\normalsize\bfseries}}
\newcommand*\l@subsubsubsection{\@dottedtocline{3}{10.0em}{4.1em}}
\makeatother
\newtheorem{theorem}{Định lý}
\newtheorem{property}{Tính chất}
\newtheorem{proposition}{Mệnh đề}
\newtheorem{corollary}[proposition]{Hệ quả}
\newtheorem{lemma}[proposition]{Bổ đề}
\theoremstyle{definition}
\newtheorem{exer}{Bài toán}
\captionsetup[figure]{labelfont={small,bf},textfont={small,it},belowskip=-1pt,aboveskip=-9pt}
\captionsetup[table]{labelfont={small,bf},textfont={small,it},belowskip=-1pt,aboveskip=7pt}
\setlength{\floatsep}{5pt plus 2pt minus 2pt}
\setlength{\textfloatsep}{5pt plus 2pt minus 2pt}
\setlength{\intextsep}{10pt plus 2pt minus 2pt}
\lstset{basicstyle=\footnotesize\sffamily,commentstyle=\ttfamily\color{black},numbers=left,numberstyle=\ttfamily\color{black}\footnotesize,stepnumber=1,numbersep=5pt,backgroundcolor=\color{white},showspaces=false,showstringspaces=false,showtabs=false,frame=single,tabsize=2,captionpos=b,breaklines=true,breakatwhitespace=false,title=\lstname,keywordstyle={}}
\hypersetup{urlcolor=blue,linkcolor=black,citecolor=black,colorlinks=true}

\begin{document}
\begin{titlepage}
\begin{tikzpicture}[remember picture,overlay]
  \draw[line width=4pt] ($(current page.north west)+(0.4in,-0.5in)$) rectangle ($(current page.south east)+(-0.4in,0.5in)$);
  \draw[line width=1.5pt] ($(current page.north west)+(0.45in,-0.55in)$) rectangle ($(current page.south east)+(-0.45in,0.55in)$);
\end{tikzpicture}
\begin{center}
{\LARGE\bfseries ĐẠI HỌC QUỐC GIA THÀNH PHỐ HỒ CHÍ MINH\par}
\vspace{0.12cm}
{\LARGE\bfseries TRƯỜNG ĐẠI HỌC BÁCH KHOA\par}
\vspace{0.12cm}
{\LARGE\bfseries KHOA KHOA HỌC VÀ KỸ THUẬT MÁY TÍNH\par}
\vspace{0.18cm}
\includegraphics[width=4cm]{Images/hcmut.png}
\vspace{0.14cm}
{\LARGE\bfseries ĐỒ ÁN MÔN HỌC\par}
\vspace{0.28cm}
{\LARGE\bfseries TRÌNH SOẠN THẢO BÁO CÁO MARKDOWN\par}
\vspace{0.28cm}
{\Large NGÀNH: KHOA HỌC MÁY TÍNH\par}
\vspace{0.25cm}
\begin{tabular}{rl}
\textbf{\Large HỘI ĐỒNG:} & {\Large KHOA HỌC MÁY TÍNH} \\[0.28cm]
\textbf{\Large GVHD:} & {\Large TS. NGUYỄN VĂN A} \\[0.28cm]
\textbf{\Large TKHĐ:} & {\Large TS. TRẦN THỊ B} \\[0.28cm]
\multicolumn{2}{c}{\textbf{\Large ---o0o---}} \\[0.28cm]
\textbf{\Large SVTH1:} & {\Large Nguyễn Văn An (2210001)} \\[0.28cm]
\textbf{\Large SVTH2:} & {\Large Trần Thị Bình (2210002)} \\[0.28cm]
\textbf{\Large SVTH3:} & {\Large Lê Minh Châu (2210003)} \\[0.28cm]
\end{tabular}
\vfill
{\Large\bfseries TP. HỒ CHÍ MINH, 2026\par}
\vspace*{1.2cm}
\end{center}
\end{titlepage}

\clearpage
\tableofcontents
\clearpage
\listoffigures
\clearpage
\listoftables
\clearpage
\setcounter{page}{1}

\section{Giới thiệu}
Tài liệu kiểm thử tiếng Việt, \textbf{chữ đậm}, \emph{chữ nghiêng} và công thức $a^2+b^2=c^2$.

\subsection{Mức hai}
\subsubsection{Mức ba}
\subsubsubsection{Mức bốn}
\begin{theorem}[Định lý mẫu]
\label{thm-golden}
Với mọi số thực $x$, ta có $x^2 \geq 0$.
\end{theorem}
\begin{property}
\label{prop-golden}
Tính chất mẫu.
\end{property}
\begin{proposition}
\label{proposition-golden}
Mệnh đề mẫu.
\end{proposition}
\begin{corollary}
\label{corollary-golden}
Hệ quả mẫu.
\end{corollary}
\begin{lemma}
\label{lemma-golden}
Bổ đề mẫu.
\end{lemma}
\begin{exer}
\label{exercise-golden}
Bài toán mẫu.
\end{exer}

\begin{table}[h!]
\centering
\begin{tabular}{ll}
\toprule
Cột A & Cột B \\
\midrule
Một & Hai \\
\bottomrule
\end{tabular}
\caption{Bảng dữ liệu mẫu}
\label{tbl-golden}
\end{table}

\begin{lstlisting}[language=swift,numbers=left,caption={Mã nguồn mẫu},label={lst-golden}]
print("Xin chào HCMUT")
\end{lstlisting}
\end{document}
